% ---------------------------------------------------------------------------
% Notation and elements of the TWFE difference-in-differences event study.
% Input into the methods note with:  % ---------------------------------------------------------------------------
% Notation and elements of the TWFE difference-in-differences event study.
% Input into the methods note with:  % ---------------------------------------------------------------------------
% Notation and elements of the TWFE difference-in-differences event study.
% Input into the methods note with:  % ---------------------------------------------------------------------------
% Notation and elements of the TWFE difference-in-differences event study.
% Input into the methods note with:  \input{tab_twfe_notation}
% Requires: booktabs, array, amsmath, longtable
%   (the methods note already loads the first three; add \usepackage{longtable}).
% Macros are guarded with \providecommand so the file also compiles standalone.
% ---------------------------------------------------------------------------
\providecommand{\Tr}{\mathcal{T}}
\providecommand{\Dn}{\mathcal{D}}
\providecommand{\ind}{\mathbf{1}}
\providecommand{\lp}{log-\%}

\begingroup
\footnotesize
\setlength{\tabcolsep}{4pt}
\renewcommand{\arraystretch}{1.12}

\begin{center}\small
Estimating equation, fitted separately for each of the six outcomes:
\[
y_{it}\;=\;\alpha_i\;+\;\gamma_t\;+\;\sum_{k\in\mathcal{K}}\beta_k\,D^{k}_{it}\;+\;\varepsilon_{it},
\qquad
D^{k}_{it}=\ind\{i\in\Tr\}\cdot\ind\{k_t=k\},
\]
by weighted least squares with analytic weights $\omega_i$, county fixed effects absorbed,
and standard errors clustered by county.
\end{center}

\begin{longtable}{@{}l>{\raggedright}p{5.4cm}>{\raggedright}p{3.9cm}>{\raggedright\arraybackslash}p{3.7cm}@{}}
\caption{Elements of the main TWFE difference-in-differences event study.}
\label{tab:twfe-notation}\\
\toprule
\textbf{Symbol} & \textbf{Meaning} & \textbf{Range / dimension} & \textbf{In the code}\\
\midrule
\endfirsthead
\multicolumn{4}{@{}l}{\textit{Table \thetable, continued from the previous page.}}\\[2pt]
\toprule
\textbf{Symbol} & \textbf{Meaning} & \textbf{Range / dimension} & \textbf{In the code}\\
\midrule
\endhead
\midrule
\multicolumn{4}{r@{}}{\textit{continued on the next page}}\\
\endfoot
\bottomrule
\endlastfoot

\multicolumn{4}{@{}l}{\textit{Panel A. Indices and estimation sample}}\\[2pt]
$i$ & County, the unit of observation. Treated counties and active donors only. & $154$ counties $=39$ treated $+\,115$ donors with $\hat w_i>10^{-10}$ & \texttt{county\_id}; selector \texttt{treated\_nc==1 | (mc\_w\_joint>1e-10 \& mc\_donor==1)}\\
$t$ & Calendar month. & $49$ months: Jan 2022--Mar 2026, \emph{less} Sep 2024 (storm transition) and Oct 2025 (LAUS missing) & \texttt{mdate}\\
$k_t$ & Event time: months since the storm month. & $k_t=t-\text{Sep\,2024}$; $49$ values in $\{-32,\dots,18\}\setminus\{0,13\}$ & \texttt{mc\_k = mdate - ym(2024,9)}\\
$\Tr$ & Frozen treated set: NC counties with positive verified FEMA loss. & $|\Tr|=39$ & \texttt{treated\_nc==1}; roster \texttt{frozen\_treated\_nc\_39.csv}\\
$\Dn$ & Donor counties entering the regression. & $115$ of the $1{,}150$ eligible donors carry positive weight & \texttt{mc\_donor==1} and $\hat w_i>10^{-10}$\\
$N$ & County-month observations. & $154\times49=7{,}546$ & \texttt{e(N)}\\[4pt]

\multicolumn{4}{@{}l}{\textit{Panel B. Dependent variable (six separate regressions)}}\\[2pt]
$y_{it}$ & Employment, labor force, workplace jobs. & natural log in the regression; reported $\times100$ as \lp & \texttt{mc\_yemp}, \texttt{mc\_ylf}, \texttt{mc\_yjobs}: raw $100\ln$ series $/100$\\
 & Unemployment rate. & percentage points, reported as is & \texttt{mc\_yurate}\\
 & Jobs\,/\,employment and employment\,/\,labor force. & log ratios: differences of the three log outcomes & \texttt{mc\_ydiff}, \texttt{mc\_yshare}\\[4pt]

\multicolumn{4}{@{}l}{\textit{Panel C. Right-hand side}}\\[2pt]
$\alpha_i$ & County fixed effect. Absorbs the permanent treated-minus-synthetic level gap that the weight-fitting intercept leaves free. & $154$ groups, $153$ free parameters & \texttt{absorb(county\_id)}\\
$\gamma_t$ & Calendar-month fixed effect: the common time path. & $49$ months, $48$ free dummies & \texttt{i.mdate}\\
$D^{k}_{it}$ & Treated $\times$ event-month indicator: $1$ for a treated county in event month $k$, else $0$. & $48$ regressors & \texttt{mc\_ev\_m32} \dots \texttt{mc\_ev\_p18}\\
$\mathcal{K}$ & Event months carrying an indicator. & $\{-32,\dots,18\}\setminus\{-1,0,13\}$: $k=-1$ (Aug 2024) is the omitted category; $k=0,13$ are not in the sample & loop \texttt{forvalues k=-32/18}, skipping $-1,0,13$\\
$\beta_k$ & Raw coefficient: treated-minus-synthetic gap in event month $k$, measured from the omitted month. Not the reported quantity. & $48$ estimated, with $\beta_{-1}\equiv0$ & \texttt{\_b[mc\_ev\_\dots]}\\
$\omega_i$ & Analytic weight, time-invariant. Treated group and donor group each carry total mass $1$. & $\omega_i=1/39$ for $i\in\Tr$; $\omega_i=\hat w_i$ for $i\in\Dn$, with $\sum_{\Dn}\hat w_i=1$ & \texttt{[aw=mc\_w\_joint]}, set to $1/39$ for treated\\
$\varepsilon_{it}$ & Error. Clustered by county: arbitrary serial correlation within a county, independence across counties. & $154$ clusters; $\mathrm{df}=153$ for every $t$ test and critical value & \texttt{vce(cluster county\_id)}; \texttt{e(df\_r)}$=153$\\[4pt]

\multicolumn{4}{@{}l}{\textit{Panel D. Reported quantities (linear combinations, joint covariance)}}\\[2pt]
$\tilde\beta_k$ & \textbf{Reported} monthly effect, recentred on the pre-event mean instead of one month: $\tilde\beta_k=\beta_k-\frac{1}{32}\sum_{m=-32}^{-1}\beta_m$, with $\beta_{-1}\equiv0$. & $49$ reported values, including $k=-1$; the pre-event $\tilde\beta_k$ average exactly $0$ & \texttt{lincom mc\_ev\_`s' - (presum)/32}\\
$\hat\tau_W$ & \textbf{Reported} window effect: the average of $\tilde\beta_k$ over window $W$. & $6$ outcomes $\times\,3$ windows $=18$ tests & \texttt{lincom (terms)/len - (presum)/32}\\
$W$ & The three pre-specified windows. & immediate $k\in\{1,2,3\}$; 2025 $k\in\{4,\dots,12\}$; late $k\in\{14,\dots,18\}$ ($k=13$ unavailable) & \texttt{immediate}, \texttt{year2025}, \texttt{late}\\
 & Multiplicity: one Holm family over the $18$ window tests. Displayed intervals are individual, not simultaneous. & Holm step-down over $18$ ordered raw $p$-values & \texttt{05\_tables\_figures.do}\\
\end{longtable}

\vspace{-6pt}
\noindent\begin{minipage}{\textwidth}\scriptsize
\textit{Notes.} Parameter count: $48$ event indicators $+\;48$ month dummies $+\;153$ absorbed county effects $+\;1$ constant $=250$, against $N=7{,}546$. Inference uses $\mathrm{df}=G-1=153$, the cluster count less one, not the residual degrees of freedom, because the variance estimator is clustered. The omitted category $k=-1$ fixes the parameterisation only: Panel D re-expresses every coefficient against the 32-month pre-event mean, so August 2024 receives its own estimate and the choice of omitted month affects no reported number. Joint Wald tests on the pre-event coefficients are invariant to this re-expression. The pre-event window holds 32 months, of which 31 carry an indicator. Implementation: \texttt{src/labor/steps/04\_joint\_event\_study.do}.
\end{minipage}
\endgroup

\begin{table}[htbp]
\centering
\caption{Additional elements of the exposure-interacted event study (exploratory).}
\label{tab:dose-notation}
\begin{minipage}{\textwidth}\centering\small
The damage dose enters the same regression as an interaction:
\[
y_{it}\;=\;\alpha_i\;+\;\gamma_t\;+\;\sum_{k\in\mathcal{K}}\Big[\beta_k\,D^{k}_{it}\;+\;\delta_k\,D^{k}_{it}\,z_i\Big]\;+\;\varepsilon_{it}.
\]
\end{minipage}
\vspace{4pt}

\footnotesize
\setlength{\tabcolsep}{4pt}
\renewcommand{\arraystretch}{1.12}
\begin{tabular}{@{}l>{\raggedright}p{5.4cm}>{\raggedright}p{3.9cm}>{\raggedright\arraybackslash}p{3.7cm}@{}}
\toprule
\textbf{Symbol} & \textbf{Meaning} & \textbf{Range / dimension} & \textbf{In the code}\\
\midrule
$z_i$ & Storm exposure, standardised over the $39$ treated counties; zero for donors, so donors never enter the interaction. & mean $0$, SD $1$ on $\Tr$; one SD $=22.4$ HDSI points or $94.7$\,mm of rain & \texttt{md\_z\_hdsi}, \texttt{md\_z\_rain}\\
$\beta_k$ & \textbf{Level} term: the effect in month $k$ for a treated county of \emph{average} exposure ($z_i=0$). & $48$ regressors & \texttt{md\_ev\_\dots}\\
$\delta_k$ & \textbf{Per-SD} term: the additional effect per one SD of exposure. This is the gradient the Gaussian curves estimate, now inside the event study. & $48$ regressors & \texttt{md\_dh\_\dots} (HDSI), \texttt{md\_dr\_\dots} (rainfall)\\
$\tilde\delta_k$ & Reported per-SD coefficient, recentred on its own pre-event mean exactly as $\tilde\beta_k$ is. & pre-event $\tilde\delta_k$ average $0$; their joint test asks whether exposure \emph{already} predicted the gap & \texttt{lincom md\_dh\_`s' - (presum)/32}\\
$g(i)$ & Tercile family: the HDSI tercile of treated county $i$, giving one event-study series per tercile instead of a level/slope split. & $3$ terciles of $13$ counties; HDSI $25.1$--$49.0$, $49.4$--$78.0$, $79.2$--$98.1$; $3\times48=144$ regressors & \texttt{md\_ter}; \texttt{md\_t1\_\dots}--\texttt{md\_t3\_\dots}\\
 & High-minus-low window contrast, from the joint covariance of the two tercile series. & $18$ contrasts & \texttt{lincom} of the two window averages\\
\bottomrule
\end{tabular}

\vspace{2pt}
\begin{minipage}{\textwidth}\scriptsize
\textit{Notes.} Sample, weights $\omega_i$, fixed effects, clustering, windows and the recentring of Panel D above are unchanged; only the right-hand side gains the interaction. Holm is applied \emph{within} each exploratory $18$-test family (per-SD HDSI; per-SD rainfall; high minus low), separately from the main family of Table~\ref{tab:twfe-notation}. A per-SD coefficient is the slope of the storm effect in exposure, not an effect level: exposure is not randomly assigned across the $39$ treated counties, and the pre-event per-SD tests reject exact month-by-month zero, so each gradient should be read against its own pre-event noise. Implementation: \texttt{src/labor/steps/04c\_dose\_event\_study.do}.
\end{minipage}
\end{table}

% Requires: booktabs, array, amsmath, longtable
%   (the methods note already loads the first three; add \usepackage{longtable}).
% Macros are guarded with \providecommand so the file also compiles standalone.
% ---------------------------------------------------------------------------
\providecommand{\Tr}{\mathcal{T}}
\providecommand{\Dn}{\mathcal{D}}
\providecommand{\ind}{\mathbf{1}}
\providecommand{\lp}{log-\%}

\begingroup
\footnotesize
\setlength{\tabcolsep}{4pt}
\renewcommand{\arraystretch}{1.12}

\begin{center}\small
Estimating equation, fitted separately for each of the six outcomes:
\[
y_{it}\;=\;\alpha_i\;+\;\gamma_t\;+\;\sum_{k\in\mathcal{K}}\beta_k\,D^{k}_{it}\;+\;\varepsilon_{it},
\qquad
D^{k}_{it}=\ind\{i\in\Tr\}\cdot\ind\{k_t=k\},
\]
by weighted least squares with analytic weights $\omega_i$, county fixed effects absorbed,
and standard errors clustered by county.
\end{center}

\begin{longtable}{@{}l>{\raggedright}p{5.4cm}>{\raggedright}p{3.9cm}>{\raggedright\arraybackslash}p{3.7cm}@{}}
\caption{Elements of the main TWFE difference-in-differences event study.}
\label{tab:twfe-notation}\\
\toprule
\textbf{Symbol} & \textbf{Meaning} & \textbf{Range / dimension} & \textbf{In the code}\\
\midrule
\endfirsthead
\multicolumn{4}{@{}l}{\textit{Table \thetable, continued from the previous page.}}\\[2pt]
\toprule
\textbf{Symbol} & \textbf{Meaning} & \textbf{Range / dimension} & \textbf{In the code}\\
\midrule
\endhead
\midrule
\multicolumn{4}{r@{}}{\textit{continued on the next page}}\\
\endfoot
\bottomrule
\endlastfoot

\multicolumn{4}{@{}l}{\textit{Panel A. Indices and estimation sample}}\\[2pt]
$i$ & County, the unit of observation. Treated counties and active donors only. & $154$ counties $=39$ treated $+\,115$ donors with $\hat w_i>10^{-10}$ & \texttt{county\_id}; selector \texttt{treated\_nc==1 | (mc\_w\_joint>1e-10 \& mc\_donor==1)}\\
$t$ & Calendar month. & $49$ months: Jan 2022--Mar 2026, \emph{less} Sep 2024 (storm transition) and Oct 2025 (LAUS missing) & \texttt{mdate}\\
$k_t$ & Event time: months since the storm month. & $k_t=t-\text{Sep\,2024}$; $49$ values in $\{-32,\dots,18\}\setminus\{0,13\}$ & \texttt{mc\_k = mdate - ym(2024,9)}\\
$\Tr$ & Frozen treated set: NC counties with positive verified FEMA loss. & $|\Tr|=39$ & \texttt{treated\_nc==1}; roster \texttt{frozen\_treated\_nc\_39.csv}\\
$\Dn$ & Donor counties entering the regression. & $115$ of the $1{,}150$ eligible donors carry positive weight & \texttt{mc\_donor==1} and $\hat w_i>10^{-10}$\\
$N$ & County-month observations. & $154\times49=7{,}546$ & \texttt{e(N)}\\[4pt]

\multicolumn{4}{@{}l}{\textit{Panel B. Dependent variable (six separate regressions)}}\\[2pt]
$y_{it}$ & Employment, labor force, workplace jobs. & natural log in the regression; reported $\times100$ as \lp & \texttt{mc\_yemp}, \texttt{mc\_ylf}, \texttt{mc\_yjobs}: raw $100\ln$ series $/100$\\
 & Unemployment rate. & percentage points, reported as is & \texttt{mc\_yurate}\\
 & Jobs\,/\,employment and employment\,/\,labor force. & log ratios: differences of the three log outcomes & \texttt{mc\_ydiff}, \texttt{mc\_yshare}\\[4pt]

\multicolumn{4}{@{}l}{\textit{Panel C. Right-hand side}}\\[2pt]
$\alpha_i$ & County fixed effect. Absorbs the permanent treated-minus-synthetic level gap that the weight-fitting intercept leaves free. & $154$ groups, $153$ free parameters & \texttt{absorb(county\_id)}\\
$\gamma_t$ & Calendar-month fixed effect: the common time path. & $49$ months, $48$ free dummies & \texttt{i.mdate}\\
$D^{k}_{it}$ & Treated $\times$ event-month indicator: $1$ for a treated county in event month $k$, else $0$. & $48$ regressors & \texttt{mc\_ev\_m32} \dots \texttt{mc\_ev\_p18}\\
$\mathcal{K}$ & Event months carrying an indicator. & $\{-32,\dots,18\}\setminus\{-1,0,13\}$: $k=-1$ (Aug 2024) is the omitted category; $k=0,13$ are not in the sample & loop \texttt{forvalues k=-32/18}, skipping $-1,0,13$\\
$\beta_k$ & Raw coefficient: treated-minus-synthetic gap in event month $k$, measured from the omitted month. Not the reported quantity. & $48$ estimated, with $\beta_{-1}\equiv0$ & \texttt{\_b[mc\_ev\_\dots]}\\
$\omega_i$ & Analytic weight, time-invariant. Treated group and donor group each carry total mass $1$. & $\omega_i=1/39$ for $i\in\Tr$; $\omega_i=\hat w_i$ for $i\in\Dn$, with $\sum_{\Dn}\hat w_i=1$ & \texttt{[aw=mc\_w\_joint]}, set to $1/39$ for treated\\
$\varepsilon_{it}$ & Error. Clustered by county: arbitrary serial correlation within a county, independence across counties. & $154$ clusters; $\mathrm{df}=153$ for every $t$ test and critical value & \texttt{vce(cluster county\_id)}; \texttt{e(df\_r)}$=153$\\[4pt]

\multicolumn{4}{@{}l}{\textit{Panel D. Reported quantities (linear combinations, joint covariance)}}\\[2pt]
$\tilde\beta_k$ & \textbf{Reported} monthly effect, recentred on the pre-event mean instead of one month: $\tilde\beta_k=\beta_k-\frac{1}{32}\sum_{m=-32}^{-1}\beta_m$, with $\beta_{-1}\equiv0$. & $49$ reported values, including $k=-1$; the pre-event $\tilde\beta_k$ average exactly $0$ & \texttt{lincom mc\_ev\_`s' - (presum)/32}\\
$\hat\tau_W$ & \textbf{Reported} window effect: the average of $\tilde\beta_k$ over window $W$. & $6$ outcomes $\times\,3$ windows $=18$ tests & \texttt{lincom (terms)/len - (presum)/32}\\
$W$ & The three pre-specified windows. & immediate $k\in\{1,2,3\}$; 2025 $k\in\{4,\dots,12\}$; late $k\in\{14,\dots,18\}$ ($k=13$ unavailable) & \texttt{immediate}, \texttt{year2025}, \texttt{late}\\
 & Multiplicity: one Holm family over the $18$ window tests. Displayed intervals are individual, not simultaneous. & Holm step-down over $18$ ordered raw $p$-values & \texttt{05\_tables\_figures.do}\\
\end{longtable}

\vspace{-6pt}
\noindent\begin{minipage}{\textwidth}\scriptsize
\textit{Notes.} Parameter count: $48$ event indicators $+\;48$ month dummies $+\;153$ absorbed county effects $+\;1$ constant $=250$, against $N=7{,}546$. Inference uses $\mathrm{df}=G-1=153$, the cluster count less one, not the residual degrees of freedom, because the variance estimator is clustered. The omitted category $k=-1$ fixes the parameterisation only: Panel D re-expresses every coefficient against the 32-month pre-event mean, so August 2024 receives its own estimate and the choice of omitted month affects no reported number. Joint Wald tests on the pre-event coefficients are invariant to this re-expression. The pre-event window holds 32 months, of which 31 carry an indicator. Implementation: \texttt{src/labor/steps/04\_joint\_event\_study.do}.
\end{minipage}
\endgroup

\begin{table}[htbp]
\centering
\caption{Additional elements of the exposure-interacted event study (exploratory).}
\label{tab:dose-notation}
\begin{minipage}{\textwidth}\centering\small
The damage dose enters the same regression as an interaction:
\[
y_{it}\;=\;\alpha_i\;+\;\gamma_t\;+\;\sum_{k\in\mathcal{K}}\Big[\beta_k\,D^{k}_{it}\;+\;\delta_k\,D^{k}_{it}\,z_i\Big]\;+\;\varepsilon_{it}.
\]
\end{minipage}
\vspace{4pt}

\footnotesize
\setlength{\tabcolsep}{4pt}
\renewcommand{\arraystretch}{1.12}
\begin{tabular}{@{}l>{\raggedright}p{5.4cm}>{\raggedright}p{3.9cm}>{\raggedright\arraybackslash}p{3.7cm}@{}}
\toprule
\textbf{Symbol} & \textbf{Meaning} & \textbf{Range / dimension} & \textbf{In the code}\\
\midrule
$z_i$ & Storm exposure, standardised over the $39$ treated counties; zero for donors, so donors never enter the interaction. & mean $0$, SD $1$ on $\Tr$; one SD $=22.4$ HDSI points or $94.7$\,mm of rain & \texttt{md\_z\_hdsi}, \texttt{md\_z\_rain}\\
$\beta_k$ & \textbf{Level} term: the effect in month $k$ for a treated county of \emph{average} exposure ($z_i=0$). & $48$ regressors & \texttt{md\_ev\_\dots}\\
$\delta_k$ & \textbf{Per-SD} term: the additional effect per one SD of exposure. This is the gradient the Gaussian curves estimate, now inside the event study. & $48$ regressors & \texttt{md\_dh\_\dots} (HDSI), \texttt{md\_dr\_\dots} (rainfall)\\
$\tilde\delta_k$ & Reported per-SD coefficient, recentred on its own pre-event mean exactly as $\tilde\beta_k$ is. & pre-event $\tilde\delta_k$ average $0$; their joint test asks whether exposure \emph{already} predicted the gap & \texttt{lincom md\_dh\_`s' - (presum)/32}\\
$g(i)$ & Tercile family: the HDSI tercile of treated county $i$, giving one event-study series per tercile instead of a level/slope split. & $3$ terciles of $13$ counties; HDSI $25.1$--$49.0$, $49.4$--$78.0$, $79.2$--$98.1$; $3\times48=144$ regressors & \texttt{md\_ter}; \texttt{md\_t1\_\dots}--\texttt{md\_t3\_\dots}\\
 & High-minus-low window contrast, from the joint covariance of the two tercile series. & $18$ contrasts & \texttt{lincom} of the two window averages\\
\bottomrule
\end{tabular}

\vspace{2pt}
\begin{minipage}{\textwidth}\scriptsize
\textit{Notes.} Sample, weights $\omega_i$, fixed effects, clustering, windows and the recentring of Panel D above are unchanged; only the right-hand side gains the interaction. Holm is applied \emph{within} each exploratory $18$-test family (per-SD HDSI; per-SD rainfall; high minus low), separately from the main family of Table~\ref{tab:twfe-notation}. A per-SD coefficient is the slope of the storm effect in exposure, not an effect level: exposure is not randomly assigned across the $39$ treated counties, and the pre-event per-SD tests reject exact month-by-month zero, so each gradient should be read against its own pre-event noise. Implementation: \texttt{src/labor/steps/04c\_dose\_event\_study.do}.
\end{minipage}
\end{table}

% Requires: booktabs, array, amsmath, longtable
%   (the methods note already loads the first three; add \usepackage{longtable}).
% Macros are guarded with \providecommand so the file also compiles standalone.
% ---------------------------------------------------------------------------
\providecommand{\Tr}{\mathcal{T}}
\providecommand{\Dn}{\mathcal{D}}
\providecommand{\ind}{\mathbf{1}}
\providecommand{\lp}{log-\%}

\begingroup
\footnotesize
\setlength{\tabcolsep}{4pt}
\renewcommand{\arraystretch}{1.12}

\begin{center}\small
Estimating equation, fitted separately for each of the six outcomes:
\[
y_{it}\;=\;\alpha_i\;+\;\gamma_t\;+\;\sum_{k\in\mathcal{K}}\beta_k\,D^{k}_{it}\;+\;\varepsilon_{it},
\qquad
D^{k}_{it}=\ind\{i\in\Tr\}\cdot\ind\{k_t=k\},
\]
by weighted least squares with analytic weights $\omega_i$, county fixed effects absorbed,
and standard errors clustered by county.
\end{center}

\begin{longtable}{@{}l>{\raggedright}p{5.4cm}>{\raggedright}p{3.9cm}>{\raggedright\arraybackslash}p{3.7cm}@{}}
\caption{Elements of the main TWFE difference-in-differences event study.}
\label{tab:twfe-notation}\\
\toprule
\textbf{Symbol} & \textbf{Meaning} & \textbf{Range / dimension} & \textbf{In the code}\\
\midrule
\endfirsthead
\multicolumn{4}{@{}l}{\textit{Table \thetable, continued from the previous page.}}\\[2pt]
\toprule
\textbf{Symbol} & \textbf{Meaning} & \textbf{Range / dimension} & \textbf{In the code}\\
\midrule
\endhead
\midrule
\multicolumn{4}{r@{}}{\textit{continued on the next page}}\\
\endfoot
\bottomrule
\endlastfoot

\multicolumn{4}{@{}l}{\textit{Panel A. Indices and estimation sample}}\\[2pt]
$i$ & County, the unit of observation. Treated counties and active donors only. & $154$ counties $=39$ treated $+\,115$ donors with $\hat w_i>10^{-10}$ & \texttt{county\_id}; selector \texttt{treated\_nc==1 | (mc\_w\_joint>1e-10 \& mc\_donor==1)}\\
$t$ & Calendar month. & $49$ months: Jan 2022--Mar 2026, \emph{less} Sep 2024 (storm transition) and Oct 2025 (LAUS missing) & \texttt{mdate}\\
$k_t$ & Event time: months since the storm month. & $k_t=t-\text{Sep\,2024}$; $49$ values in $\{-32,\dots,18\}\setminus\{0,13\}$ & \texttt{mc\_k = mdate - ym(2024,9)}\\
$\Tr$ & Frozen treated set: NC counties with positive verified FEMA loss. & $|\Tr|=39$ & \texttt{treated\_nc==1}; roster \texttt{frozen\_treated\_nc\_39.csv}\\
$\Dn$ & Donor counties entering the regression. & $115$ of the $1{,}150$ eligible donors carry positive weight & \texttt{mc\_donor==1} and $\hat w_i>10^{-10}$\\
$N$ & County-month observations. & $154\times49=7{,}546$ & \texttt{e(N)}\\[4pt]

\multicolumn{4}{@{}l}{\textit{Panel B. Dependent variable (six separate regressions)}}\\[2pt]
$y_{it}$ & Employment, labor force, workplace jobs. & natural log in the regression; reported $\times100$ as \lp & \texttt{mc\_yemp}, \texttt{mc\_ylf}, \texttt{mc\_yjobs}: raw $100\ln$ series $/100$\\
 & Unemployment rate. & percentage points, reported as is & \texttt{mc\_yurate}\\
 & Jobs\,/\,employment and employment\,/\,labor force. & log ratios: differences of the three log outcomes & \texttt{mc\_ydiff}, \texttt{mc\_yshare}\\[4pt]

\multicolumn{4}{@{}l}{\textit{Panel C. Right-hand side}}\\[2pt]
$\alpha_i$ & County fixed effect. Absorbs the permanent treated-minus-synthetic level gap that the weight-fitting intercept leaves free. & $154$ groups, $153$ free parameters & \texttt{absorb(county\_id)}\\
$\gamma_t$ & Calendar-month fixed effect: the common time path. & $49$ months, $48$ free dummies & \texttt{i.mdate}\\
$D^{k}_{it}$ & Treated $\times$ event-month indicator: $1$ for a treated county in event month $k$, else $0$. & $48$ regressors & \texttt{mc\_ev\_m32} \dots \texttt{mc\_ev\_p18}\\
$\mathcal{K}$ & Event months carrying an indicator. & $\{-32,\dots,18\}\setminus\{-1,0,13\}$: $k=-1$ (Aug 2024) is the omitted category; $k=0,13$ are not in the sample & loop \texttt{forvalues k=-32/18}, skipping $-1,0,13$\\
$\beta_k$ & Raw coefficient: treated-minus-synthetic gap in event month $k$, measured from the omitted month. Not the reported quantity. & $48$ estimated, with $\beta_{-1}\equiv0$ & \texttt{\_b[mc\_ev\_\dots]}\\
$\omega_i$ & Analytic weight, time-invariant. Treated group and donor group each carry total mass $1$. & $\omega_i=1/39$ for $i\in\Tr$; $\omega_i=\hat w_i$ for $i\in\Dn$, with $\sum_{\Dn}\hat w_i=1$ & \texttt{[aw=mc\_w\_joint]}, set to $1/39$ for treated\\
$\varepsilon_{it}$ & Error. Clustered by county: arbitrary serial correlation within a county, independence across counties. & $154$ clusters; $\mathrm{df}=153$ for every $t$ test and critical value & \texttt{vce(cluster county\_id)}; \texttt{e(df\_r)}$=153$\\[4pt]

\multicolumn{4}{@{}l}{\textit{Panel D. Reported quantities (linear combinations, joint covariance)}}\\[2pt]
$\tilde\beta_k$ & \textbf{Reported} monthly effect, recentred on the pre-event mean instead of one month: $\tilde\beta_k=\beta_k-\frac{1}{32}\sum_{m=-32}^{-1}\beta_m$, with $\beta_{-1}\equiv0$. & $49$ reported values, including $k=-1$; the pre-event $\tilde\beta_k$ average exactly $0$ & \texttt{lincom mc\_ev\_`s' - (presum)/32}\\
$\hat\tau_W$ & \textbf{Reported} window effect: the average of $\tilde\beta_k$ over window $W$. & $6$ outcomes $\times\,3$ windows $=18$ tests & \texttt{lincom (terms)/len - (presum)/32}\\
$W$ & The three pre-specified windows. & immediate $k\in\{1,2,3\}$; 2025 $k\in\{4,\dots,12\}$; late $k\in\{14,\dots,18\}$ ($k=13$ unavailable) & \texttt{immediate}, \texttt{year2025}, \texttt{late}\\
 & Multiplicity: one Holm family over the $18$ window tests. Displayed intervals are individual, not simultaneous. & Holm step-down over $18$ ordered raw $p$-values & \texttt{05\_tables\_figures.do}\\
\end{longtable}

\vspace{-6pt}
\noindent\begin{minipage}{\textwidth}\scriptsize
\textit{Notes.} Parameter count: $48$ event indicators $+\;48$ month dummies $+\;153$ absorbed county effects $+\;1$ constant $=250$, against $N=7{,}546$. Inference uses $\mathrm{df}=G-1=153$, the cluster count less one, not the residual degrees of freedom, because the variance estimator is clustered. The omitted category $k=-1$ fixes the parameterisation only: Panel D re-expresses every coefficient against the 32-month pre-event mean, so August 2024 receives its own estimate and the choice of omitted month affects no reported number. Joint Wald tests on the pre-event coefficients are invariant to this re-expression. The pre-event window holds 32 months, of which 31 carry an indicator. Implementation: \texttt{src/labor/steps/04\_joint\_event\_study.do}.
\end{minipage}
\endgroup

\begin{table}[htbp]
\centering
\caption{Additional elements of the exposure-interacted event study (exploratory).}
\label{tab:dose-notation}
\begin{minipage}{\textwidth}\centering\small
The damage dose enters the same regression as an interaction:
\[
y_{it}\;=\;\alpha_i\;+\;\gamma_t\;+\;\sum_{k\in\mathcal{K}}\Big[\beta_k\,D^{k}_{it}\;+\;\delta_k\,D^{k}_{it}\,z_i\Big]\;+\;\varepsilon_{it}.
\]
\end{minipage}
\vspace{4pt}

\footnotesize
\setlength{\tabcolsep}{4pt}
\renewcommand{\arraystretch}{1.12}
\begin{tabular}{@{}l>{\raggedright}p{5.4cm}>{\raggedright}p{3.9cm}>{\raggedright\arraybackslash}p{3.7cm}@{}}
\toprule
\textbf{Symbol} & \textbf{Meaning} & \textbf{Range / dimension} & \textbf{In the code}\\
\midrule
$z_i$ & Storm exposure, standardised over the $39$ treated counties; zero for donors, so donors never enter the interaction. & mean $0$, SD $1$ on $\Tr$; one SD $=22.4$ HDSI points or $94.7$\,mm of rain & \texttt{md\_z\_hdsi}, \texttt{md\_z\_rain}\\
$\beta_k$ & \textbf{Level} term: the effect in month $k$ for a treated county of \emph{average} exposure ($z_i=0$). & $48$ regressors & \texttt{md\_ev\_\dots}\\
$\delta_k$ & \textbf{Per-SD} term: the additional effect per one SD of exposure. This is the gradient the Gaussian curves estimate, now inside the event study. & $48$ regressors & \texttt{md\_dh\_\dots} (HDSI), \texttt{md\_dr\_\dots} (rainfall)\\
$\tilde\delta_k$ & Reported per-SD coefficient, recentred on its own pre-event mean exactly as $\tilde\beta_k$ is. & pre-event $\tilde\delta_k$ average $0$; their joint test asks whether exposure \emph{already} predicted the gap & \texttt{lincom md\_dh\_`s' - (presum)/32}\\
$g(i)$ & Tercile family: the HDSI tercile of treated county $i$, giving one event-study series per tercile instead of a level/slope split. & $3$ terciles of $13$ counties; HDSI $25.1$--$49.0$, $49.4$--$78.0$, $79.2$--$98.1$; $3\times48=144$ regressors & \texttt{md\_ter}; \texttt{md\_t1\_\dots}--\texttt{md\_t3\_\dots}\\
 & High-minus-low window contrast, from the joint covariance of the two tercile series. & $18$ contrasts & \texttt{lincom} of the two window averages\\
\bottomrule
\end{tabular}

\vspace{2pt}
\begin{minipage}{\textwidth}\scriptsize
\textit{Notes.} Sample, weights $\omega_i$, fixed effects, clustering, windows and the recentring of Panel D above are unchanged; only the right-hand side gains the interaction. Holm is applied \emph{within} each exploratory $18$-test family (per-SD HDSI; per-SD rainfall; high minus low), separately from the main family of Table~\ref{tab:twfe-notation}. A per-SD coefficient is the slope of the storm effect in exposure, not an effect level: exposure is not randomly assigned across the $39$ treated counties, and the pre-event per-SD tests reject exact month-by-month zero, so each gradient should be read against its own pre-event noise. Implementation: \texttt{src/labor/steps/04c\_dose\_event\_study.do}.
\end{minipage}
\end{table}

% Requires: booktabs, array, amsmath, longtable
%   (the methods note already loads the first three; add \usepackage{longtable}).
% Macros are guarded with \providecommand so the file also compiles standalone.
% ---------------------------------------------------------------------------
\providecommand{\Tr}{\mathcal{T}}
\providecommand{\Dn}{\mathcal{D}}
\providecommand{\ind}{\mathbf{1}}
\providecommand{\lp}{log-\%}

\begingroup
\footnotesize
\setlength{\tabcolsep}{4pt}
\renewcommand{\arraystretch}{1.12}

\begin{center}\small
Estimating equation, fitted separately for each of the six outcomes:
\[
y_{it}\;=\;\alpha_i\;+\;\gamma_t\;+\;\sum_{k\in\mathcal{K}}\beta_k\,D^{k}_{it}\;+\;\varepsilon_{it},
\qquad
D^{k}_{it}=\ind\{i\in\Tr\}\cdot\ind\{k_t=k\},
\]
by weighted least squares with analytic weights $\omega_i$, county fixed effects absorbed,
and standard errors clustered by county.
\end{center}

\begin{longtable}{@{}l>{\raggedright}p{5.4cm}>{\raggedright}p{3.9cm}>{\raggedright\arraybackslash}p{3.7cm}@{}}
\caption{Elements of the main TWFE difference-in-differences event study.}
\label{tab:twfe-notation}\\
\toprule
\textbf{Symbol} & \textbf{Meaning} & \textbf{Range / dimension} & \textbf{In the code}\\
\midrule
\endfirsthead
\multicolumn{4}{@{}l}{\textit{Table \thetable, continued from the previous page.}}\\[2pt]
\toprule
\textbf{Symbol} & \textbf{Meaning} & \textbf{Range / dimension} & \textbf{In the code}\\
\midrule
\endhead
\midrule
\multicolumn{4}{r@{}}{\textit{continued on the next page}}\\
\endfoot
\bottomrule
\endlastfoot

\multicolumn{4}{@{}l}{\textit{Panel A. Indices and estimation sample}}\\[2pt]
$i$ & County, the unit of observation. Treated counties and active donors only. & $154$ counties $=39$ treated $+\,115$ donors with $\hat w_i>10^{-10}$ & \texttt{county\_id}; selector \texttt{treated\_nc==1 | (mc\_w\_joint>1e-10 \& mc\_donor==1)}\\
$t$ & Calendar month. & $49$ months: Jan 2022--Mar 2026, \emph{less} Sep 2024 (storm transition) and Oct 2025 (LAUS missing) & \texttt{mdate}\\
$k_t$ & Event time: months since the storm month. & $k_t=t-\text{Sep\,2024}$; $49$ values in $\{-32,\dots,18\}\setminus\{0,13\}$ & \texttt{mc\_k = mdate - ym(2024,9)}\\
$\Tr$ & Frozen treated set: NC counties with positive verified FEMA loss. & $|\Tr|=39$ & \texttt{treated\_nc==1}; roster \texttt{frozen\_treated\_nc\_39.csv}\\
$\Dn$ & Donor counties entering the regression. & $115$ of the $1{,}150$ eligible donors carry positive weight & \texttt{mc\_donor==1} and $\hat w_i>10^{-10}$\\
$N$ & County-month observations. & $154\times49=7{,}546$ & \texttt{e(N)}\\[4pt]

\multicolumn{4}{@{}l}{\textit{Panel B. Dependent variable (six separate regressions)}}\\[2pt]
$y_{it}$ & Employment, labor force, workplace jobs. & natural log in the regression; reported $\times100$ as \lp & \texttt{mc\_yemp}, \texttt{mc\_ylf}, \texttt{mc\_yjobs}: raw $100\ln$ series $/100$\\
 & Unemployment rate. & percentage points, reported as is & \texttt{mc\_yurate}\\
 & Jobs\,/\,employment and employment\,/\,labor force. & log ratios: differences of the three log outcomes & \texttt{mc\_ydiff}, \texttt{mc\_yshare}\\[4pt]

\multicolumn{4}{@{}l}{\textit{Panel C. Right-hand side}}\\[2pt]
$\alpha_i$ & County fixed effect. Absorbs the permanent treated-minus-synthetic level gap that the weight-fitting intercept leaves free. & $154$ groups, $153$ free parameters & \texttt{absorb(county\_id)}\\
$\gamma_t$ & Calendar-month fixed effect: the common time path. & $49$ months, $48$ free dummies & \texttt{i.mdate}\\
$D^{k}_{it}$ & Treated $\times$ event-month indicator: $1$ for a treated county in event month $k$, else $0$. & $48$ regressors & \texttt{mc\_ev\_m32} \dots \texttt{mc\_ev\_p18}\\
$\mathcal{K}$ & Event months carrying an indicator. & $\{-32,\dots,18\}\setminus\{-1,0,13\}$: $k=-1$ (Aug 2024) is the omitted category; $k=0,13$ are not in the sample & loop \texttt{forvalues k=-32/18}, skipping $-1,0,13$\\
$\beta_k$ & Raw coefficient: treated-minus-synthetic gap in event month $k$, measured from the omitted month. Not the reported quantity. & $48$ estimated, with $\beta_{-1}\equiv0$ & \texttt{\_b[mc\_ev\_\dots]}\\
$\omega_i$ & Analytic weight, time-invariant. Treated group and donor group each carry total mass $1$. & $\omega_i=1/39$ for $i\in\Tr$; $\omega_i=\hat w_i$ for $i\in\Dn$, with $\sum_{\Dn}\hat w_i=1$ & \texttt{[aw=mc\_w\_joint]}, set to $1/39$ for treated\\
$\varepsilon_{it}$ & Error. Clustered by county: arbitrary serial correlation within a county, independence across counties. & $154$ clusters; $\mathrm{df}=153$ for every $t$ test and critical value & \texttt{vce(cluster county\_id)}; \texttt{e(df\_r)}$=153$\\[4pt]

\multicolumn{4}{@{}l}{\textit{Panel D. Reported quantities (linear combinations, joint covariance)}}\\[2pt]
$\tilde\beta_k$ & \textbf{Reported} monthly effect, recentred on the pre-event mean instead of one month: $\tilde\beta_k=\beta_k-\frac{1}{32}\sum_{m=-32}^{-1}\beta_m$, with $\beta_{-1}\equiv0$. & $49$ reported values, including $k=-1$; the pre-event $\tilde\beta_k$ average exactly $0$ & \texttt{lincom mc\_ev\_`s' - (presum)/32}\\
$\hat\tau_W$ & \textbf{Reported} window effect: the average of $\tilde\beta_k$ over window $W$. & $6$ outcomes $\times\,3$ windows $=18$ tests & \texttt{lincom (terms)/len - (presum)/32}\\
$W$ & The three pre-specified windows. & immediate $k\in\{1,2,3\}$; 2025 $k\in\{4,\dots,12\}$; late $k\in\{14,\dots,18\}$ ($k=13$ unavailable) & \texttt{immediate}, \texttt{year2025}, \texttt{late}\\
 & Multiplicity: one Holm family over the $18$ window tests. Displayed intervals are individual, not simultaneous. & Holm step-down over $18$ ordered raw $p$-values & \texttt{05\_tables\_figures.do}\\
\end{longtable}

\vspace{-6pt}
\noindent\begin{minipage}{\textwidth}\scriptsize
\textit{Notes.} Parameter count: $48$ event indicators $+\;48$ month dummies $+\;153$ absorbed county effects $+\;1$ constant $=250$, against $N=7{,}546$. Inference uses $\mathrm{df}=G-1=153$, the cluster count less one, not the residual degrees of freedom, because the variance estimator is clustered. The omitted category $k=-1$ fixes the parameterisation only: Panel D re-expresses every coefficient against the 32-month pre-event mean, so August 2024 receives its own estimate and the choice of omitted month affects no reported number. Joint Wald tests on the pre-event coefficients are invariant to this re-expression. The pre-event window holds 32 months, of which 31 carry an indicator. Implementation: \texttt{src/labor/steps/04\_joint\_event\_study.do}.
\end{minipage}
\endgroup

\begin{table}[htbp]
\centering
\caption{Additional elements of the exposure-interacted event study (exploratory).}
\label{tab:dose-notation}
\begin{minipage}{\textwidth}\centering\small
The damage dose enters the same regression as an interaction:
\[
y_{it}\;=\;\alpha_i\;+\;\gamma_t\;+\;\sum_{k\in\mathcal{K}}\Big[\beta_k\,D^{k}_{it}\;+\;\delta_k\,D^{k}_{it}\,z_i\Big]\;+\;\varepsilon_{it}.
\]
\end{minipage}
\vspace{4pt}

\footnotesize
\setlength{\tabcolsep}{4pt}
\renewcommand{\arraystretch}{1.12}
\begin{tabular}{@{}l>{\raggedright}p{5.4cm}>{\raggedright}p{3.9cm}>{\raggedright\arraybackslash}p{3.7cm}@{}}
\toprule
\textbf{Symbol} & \textbf{Meaning} & \textbf{Range / dimension} & \textbf{In the code}\\
\midrule
$z_i$ & Storm exposure, standardised over the $39$ treated counties; zero for donors, so donors never enter the interaction. & mean $0$, SD $1$ on $\Tr$; one SD $=22.4$ HDSI points or $94.7$\,mm of rain & \texttt{md\_z\_hdsi}, \texttt{md\_z\_rain}\\
$\beta_k$ & \textbf{Level} term: the effect in month $k$ for a treated county of \emph{average} exposure ($z_i=0$). & $48$ regressors & \texttt{md\_ev\_\dots}\\
$\delta_k$ & \textbf{Per-SD} term: the additional effect per one SD of exposure. This is the gradient the Gaussian curves estimate, now inside the event study. & $48$ regressors & \texttt{md\_dh\_\dots} (HDSI), \texttt{md\_dr\_\dots} (rainfall)\\
$\tilde\delta_k$ & Reported per-SD coefficient, recentred on its own pre-event mean exactly as $\tilde\beta_k$ is. & pre-event $\tilde\delta_k$ average $0$; their joint test asks whether exposure \emph{already} predicted the gap & \texttt{lincom md\_dh\_`s' - (presum)/32}\\
$g(i)$ & Tercile family: the HDSI tercile of treated county $i$, giving one event-study series per tercile instead of a level/slope split. & $3$ terciles of $13$ counties; HDSI $25.1$--$49.0$, $49.4$--$78.0$, $79.2$--$98.1$; $3\times48=144$ regressors & \texttt{md\_ter}; \texttt{md\_t1\_\dots}--\texttt{md\_t3\_\dots}\\
 & High-minus-low window contrast, from the joint covariance of the two tercile series. & $18$ contrasts & \texttt{lincom} of the two window averages\\
\bottomrule
\end{tabular}

\vspace{2pt}
\begin{minipage}{\textwidth}\scriptsize
\textit{Notes.} Sample, weights $\omega_i$, fixed effects, clustering, windows and the recentring of Panel D above are unchanged; only the right-hand side gains the interaction. Holm is applied \emph{within} each exploratory $18$-test family (per-SD HDSI; per-SD rainfall; high minus low), separately from the main family of Table~\ref{tab:twfe-notation}. A per-SD coefficient is the slope of the storm effect in exposure, not an effect level: exposure is not randomly assigned across the $39$ treated counties, and the pre-event per-SD tests reject exact month-by-month zero, so each gradient should be read against its own pre-event noise. Implementation: \texttt{src/labor/steps/04c\_dose\_event\_study.do}.
\end{minipage}
\end{table}
